% Created 2018-09-07 Fri 15:25
% Intended LaTeX compiler: pdflatex
\documentclass[11pt]{article}
\usepackage[utf8]{inputenc}
\usepackage[T1]{fontenc}
\usepackage{graphicx}
\usepackage{grffile}
\usepackage{longtable}
\usepackage{wrapfig}
\usepackage{rotating}
\usepackage[normalem]{ulem}
\usepackage{amsmath}
\usepackage{textcomp}
\usepackage{amssymb}
\usepackage{capt-of}
\usepackage{hyperref}
\usepackage[margin=0.80in]{geometry}
\usepackage[backend=biber, style=ieee, url=false]{biblatex}
\usepackage{float}
\usepackage[super,negative]{nth}
\usepackage[capitalise]{cleveref}
\usepackage{pst-node,transparent,ragged2e}
\addbibresource{/home/nuk3/.spacemacs.d/org-files/bibliography.bib}
\DeclareFieldFormat[inproceedings]{citetitle}{\textit{#1}}
\DeclareFieldFormat[inproceedings]{title}{\textit{#1}}
\DeclareFieldFormat[misc]{citetitle}{#1}
\DeclareFieldFormat[misc]{title}{#1}
\renewcommand*{\bibpagespunct}{%
\ifentrytype{inproceedings}
{\addspace}
{\addcomma\space}}
\AtEveryCitekey{\ifuseauthor{}{\clearname{author}}}
\AtEveryBibitem{\ifuseauthor{}{\clearname{author}}}
\SpecialCoor
\author{Nyameko Lisa}
\date{31 August 2018}
\title{IPSP089 - Legal Aspects of Traditional Knowledge and Bioduversity\\\medskip
\large Assignment 1: 781450}
\hypersetup{
 pdfauthor={Nyameko Lisa},
 pdftitle={IPSP089 - Legal Aspects of Traditional Knowledge and Bioduversity},
 pdfkeywords={},
 pdfsubject={},
 pdfcreator={Emacs 27.0.50 (Org mode 9.1.2)},
 pdflang={English}}
\begin{document}

\maketitle
\addvspace{110pt}
\centering{
\pnode(0.5\textwidth,-0.5\textheight){thisCenter}
\rput(thisCenter){%\transparent{0.25}
\includegraphics[width=2.7in]{/home/nuk3/course/llb/wipo-unisa/UNISACoatofArms.eps}}}

\justifying
\addvspace{110pt}
\section*{Declaration}
\label{sec:org5032493}
I know that plagiarism is to use someone else’s work and pass it off as my own.\\
I know that plagiarism is wrong.\\
I confirm that this assignment is my own work.\\
I have acknowledged in the bibliography accompanying the assignment all the sources that I have used.\\
I have not directly copied without acknowledgement anything from the Internet or from any other source.\\
I have indicated every quotation and citation in a footnote or bracket linked to that quotation.\\
I have not allowed anyone else to copy my work and to pass it off as their own work.\\
I understand that if any unacknowledged copying whatsoever appears in my assignment I will receive zero per cent for the assignment.\\
I am aware of the UNISA policy on plagiarism and understand that disciplinary proceedings can be instituted against me by UNISA if I contravene this policy.\\
I indicate my understanding and acceptance of this declaration by
entering my name here under:
\begin{itemize}
\item Name: \textbf{Nyameko Lisa} (Student Number: \textbf{7874-909-3})
\end{itemize}

\subsection*{NOTE}
\label{sec:org7ddb573}
Please note that footnotes will be denoted as \footnote{This is a footnote.} and will
appear at the bottom of the page.\\
References will be denoted by \cite{rsa93_tm_act} and will appear at the end of the document.
\newpage

\section{Question 1 [50]}
\label{sec:org1cc554b}

\cite{simelane09_african_trad_knowl_sys_biodiv_manag,saurombe09_protection_indigenous_trad_knowl_through_ip,masango14_trad_knowl_cultural_express_protect,masango10_indigenous_trad_knowl_protect_prospects,amechi15_trad_knowl_medical_uses_plants_pat_sa,amechi15_leverag_trad_knowl_medicinal,geyer10_towards_clearer_defin_underst_indi,rsa78_patents_act,wto17_trips}
\subsection{Compare and contrast the terms ``traditional knowledge'' and ``indigenous knowledge''. In your answer state whether the \emph{maca} plant falls within these definitions. (10)}
\label{sec:orgbad1248}
Traditional knowledge systems, creations, innovations, or cultural expressions
refer to \emph{literary, artistic or scientific works, performances, scientific
inventions or discoveries, designs, marks or symbols}, that have evolved in
response to a changing environment and are inherited from one generation to
another within a particular group of people or their territory.\\

Although there exists no universally accepted formal definition for indigenous
peoples, `indigenous knowledge' is synonymous with and in fact a subset to,
`traditional knowledge' in terms of its scope, transmission and diversity.
Whilst their distinction is indeed subtle, `traditional knowledge' held by
indigenous people, denotes `indigenous knowledge'.\\

For example, knowledge disseminated through traditional means amongst the
Khoisan peoples of Southern Africa, constitutes `indigenous' or `traditional
knowledge'. In contrast however, the information passed down through traditional
means, by colonial Dutch settlers in the Cape can be described as `traditional
knowledge', but \textbf{not} 'indigenous knowledge'.\\

Given that the maca has long been valued for its ability to enhance fertility in
humans and livestock, it does indeed constitute `traditional knowledge'.
Moreover, given that indigenous peoples have used it the same way for
generations, as a potato-like food source, as well as a supplement for
livestock, appreciated not only for its rich protein and vitamin content, but
for it's sexual performance and fertility enhancement, such use also constitutes
`indigenous knowledge'.

\subsection{How has knowledge been protected customarily? (10)}
\label{sec:org485676c}


\subsection{What are the problems confronting rights holders (in our case the people in Junin Peru)? (10)}
\label{sec:org0767d66}
But intellectual property is not only about property. It is also
about recognition of identifiable, human creators and respect
for their contributions. From this perspective, intellectual
property has a very important role to play in protecting the
dignity of the holders of traditional knowledge. By conferring
property rights in relation to such knowledge, those holders
enabled, to a certain extent, to control its use by others.


\subsection{What is WIPO's role in protecting traditional knowledge? (20)}
\label{sec:org5d9cd52}

\printbibliography
\end{document}